\documentclass[12pt]{article}
\usepackage[T1]{fontenc}
\usepackage[utf8]{inputenc}
\usepackage{lmodern}
\usepackage{textcomp}
\usepackage{enumitem}
\usepackage{geometry}
\geometry{margin=1in}
\begin{document}
\begin{center}
Answer Key - Political Science - Undergraduate Level Assessment 21 \
\small (Answers are ordered exactly as in the question paper.)
\end{center}

\section*{Multiple Choice}
\begin{enumerate}[label=\arabic*.]
\item \textbf{Answer:} A
\\ \textit{Options:} A) refuses to hear the appeal; B) declares a mistrial; C) orders the lower court to retry the case; D) rules in favor of the defendant
\\ \textit{Note:} The Supreme Court typically refuses to hear appeals because it receives thousands of requests each year but only selects a small number of cases that have significant legal or constitutional implications, prioritizing issues that have broader national importance.
\item \textbf{Answer:} A
\\ \textit{Options:} A) power is shared between the national government and state governments; B) states have equal representation in the national government; C) individual liberties are guaranteed by a Bill of Rights; D) legislative, executive, and judicial powers are separated
\\ \textit{Note:} The correct answer is accurate because federalism specifically defines a system of governance in which sovereignty is constitutionally divided between a central authority and smaller political units, allowing both the national government and state governments to exercise power independently within their respective spheres.
\item \textbf{Answer:} D
\\ \textit{Options:} A) They arose and succeeded during times of prosperity.; B) They were created to protest wars.; C) At some point they won a majority of seats in Congress.; D) They flourished during periods of widespread dissatisfaction.
\\ \textit{Note:} Most third parties in U.S. history emerged and gained traction during times of significant public discontent with the major political parties, reflecting the electorate's desire for alternative solutions and representation.
\item \textbf{Answer:} A
\\ \textit{Options:} A) Members of Congress negotiating directly with foreign governments; B) Face-to-face meetings between state leaders; C) The president consulting Congress on foreign policy issues; D) Bilateral talks that do not involve a third-party negotiator
\\ \textit{Note:} Direct diplomacy refers to the process by which members of Congress engage in negotiations and discussions directly with foreign governments, thereby influencing foreign policy and international relations without intermediaries.
\item \textbf{Answer:} D
\\ \textit{Options:} A) Terrorists; B) Organized crime; C) Drug traffickers; D) China
\\ \textit{Note:} China is considered a state actor and a major world power, whereas nonstate actors typically include organizations or groups that operate independently of government authority and may pose threats to national security.
\end{enumerate}

\section*{True / False}
\begin{enumerate}[label=\arabic*.]
\setcounter{enumi}{5}
\item \textbf{Answer:} False
\\ \textit{Options:} A) True; B) False
\\ \textit{Note:} The concept of 'realism' in international relations predates John Foster Dulles and was developed by scholars such as Hans Morgenthau, focusing on power politics and the anarchic nature of international systems, rather than being a term specifically coined by Dulles for U.S.-Soviet relations.
\item \textbf{Answer:} True
\\ \textit{Options:} A) True; B) False
\\ \textit{Note:} An unfunded mandate is defined as a regulation or policy imposed by the federal government that obligates state and local governments to perform certain actions without allocating the necessary financial resources to support those actions.
\item \textbf{Answer:} False
\\ \textit{Options:} A) True; B) False
\\ \textit{Note:} States have moved their presidential primary dates earlier primarily to gain greater influence in the nomination process, rather than to reduce the overall cost of conducting elections.
\item \textbf{Answer:} True
\\ \textit{Options:} A) True; B) False
\\ \textit{Note:} Nations often consist of shared cultural, linguistic, and historical traits that transcend political borders, allowing for the existence of multiple nations within a single state or a single nation across multiple states.
\item \textbf{Answer:} False
\\ \textit{Options:} A) True; B) False
\\ \textit{Note:} The Supreme Court case Regents of University of California v. Bakke primarily dealt with affirmative action and the legality of race-based admissions policies, rather than the rights of students to protest on campus.
\end{enumerate}

\section*{Long Form Response}
\begin{enumerate}[label=\arabic*.]
\setcounter{enumi}{10}
\item \textbf{Answer:} International institutions have been established to address a variety of challenges, including economic stability, security threats, environmental issues, and human rights violations. These institutions, such as the United Nations, World Trade Organization, and International Monetary Fund, aim to foster global cooperation and stability by providing platforms for dialogue, conflict resolution, and coordinated action. Their effectiveness, however, can be mixed; while they have successfully facilitated negotiations and established norms, they often struggle with enforcement and the varying commitment levels of member states. Additionally, political and cultural differences can hinder consensus- building, affecting their ability to respond swiftly to crises. Overall, while international institutions have made significant strides in promoting cooperation, their effectiveness is contingent upon the willingness of states to collaborate and adhere to established agreements.
\\ \textit{Note:} correct because it identifies specific challenges addressed by international institutions and evaluates their mixed effectiveness in promoting global cooperation and stability through dialogue and coordinated action.
\item \textbf{Answer:} Acknowledging the limits of absolute security in cyber defense is crucial because it fosters a realistic understanding of vulnerabilities and threats that can never be entirely eliminated. This recognition encourages the design of resilient IT systems that can withstand and recover from attacks, rather than striving for an unattainable perfect security. Additionally, international cooperation plays a vital role in enhancing security, as it facilitates information exchange and collective responses to cyber threats. By sharing intelligence and resources, nations can build a more robust defense network that mitigates risks and enhances overall resilience against potential cyber incidents. Ultimately, a combination of resilient system design and collaborative international efforts can significantly reduce threats and improve the security landscape.
\\ \textit{Note:} correct because it emphasizes the necessity of accepting the inherent vulnerabilities in cyber defense, which leads to the development of resilient systems and underscores the importance of international collaboration in enhancing overall security through shared knowledge and coordinated responses.
\end{enumerate}

\vfill
\noindent\textit{End of Answer Key}
\end{document}